\section{Related Work}
\label{sec:related_work}

This research intersects four active areas of computational linguistics and machine learning: AI-generated text detection, low-resource natural language processing for agglutinative Turkic languages, automated fact verification pipelines, and multi-dimensional trust calibration.

\subsection{Machine-Generated Text Detection}
The task of detecting text authored by Large Language Models has evolved along two primary methodological trajectories: zero-shot statistical estimators and supervised discriminative classifiers.

\paragraph{Statistical and Curvature-Based Detectors}
Early statistical approaches operated on the hypothesis that autoregressive language models sample words disproportionately from regions of high conditional probability compared to human authors. The Giant Language Model Test Room (GLTR)~\citep{gehrmann2019gltr} introduced visual inspection and statistical hypothesis testing over token rank distributions, demonstrating that machine text clusters heavily within the top-$k$ predictive ranks. DetectGPT~\citep{mitchell2023detectgpt} advanced this paradigm by formalizing probability curvature: evaluating the log-probability differential between a candidate text and its locally perturbed variants under a surrogate language model. Fast-DetectGPT~\citep{bao2024fastdetectgpt} subsequently accelerated this process by substituting perturbation sampling with conditional probability curvature, while Binoculars~\citep{hans2024binoculars} leveraged cross-perplexity metrics between paired base and chat models to achieve zero-shot detection without explicit training. 

Concurrently, generation-phase watermarking algorithms~\citep{kirchenbauer2023watermark} embed subtle pseudo-random signals directly into LLM decoding logits by partitioning vocabulary tokens into green and red lists. While mathematically elegant, watermarking approaches necessitate white-box access to model generation APIs, rendering them inapplicable for auditing arbitrary third-party texts circulating on public digital platforms.

\paragraph{Supervised Neural Classifiers and Generalization Vulnerabilities}
Supervised fine-tuning of deep bidirectional transformers (such as BERT and RoBERTa) on balanced pairs of human and machine texts remains the most empirically performant detection methodology in practice~\citep{solaiman2019release,uchendu2020authorship,ippolito2020automatic}. However, recent benchmark evaluations have exposed severe vulnerabilities in supervised detectors. Sadasivan et al.~\citep{sadasivan2023can} and Chakraborty et al.~\citep{chakraborty2023possibilities} demonstrated that neural detectors are susceptible to paraphrase attacks and decay rapidly when tested across generation prompts, decoding strategies, or model architectures outside their training distribution. 

The comprehensive RAID benchmark~\citep{wang2024raid} highlighted that detectors trained on specific genres frequently overfit to superficial stylistic artifacts rather than true generative signals. Crucially, Liang et al.~\citep{liang2023gptdetectors} revealed that statistical AI detectors exhibit acute biases against non-native English speakers, mistakenly classifying authentic human writing as machine-generated due to constrained vocabulary usage and lower perplexity. In this work, we demonstrate that a parallel failure mode afflicts agglutinative languages: standard subword tokenizers partition native morphology into erratic fragments, inducing catastrophic false positive rates on short human texts.

\subsection{Low-Resource and Turkic NLP: The Agglutinative Challenge}
Turkic languages, encompassing Kazakh, Turkish, Uzbek, Kyrgyz, and Uyghur, are characterized by highly regular, concatenative agglutinative morphotactics~\citep{oflazer1994twolevel,bekmanova2021kazakh}. Suffixes attach in strictly governed linear tiers to signify grammatical cases, possessives, verbal tenses, aspects, evidentiality, and negation.

\paragraph{Morphological Analysis and Finite-State Transducers}
Computational processing of agglutinative languages has historically relied on two-level morphology and Finite-State Transducers (FST). Oflazer~\citep{oflazer1994twolevel} pioneered two-level morphological analyzers for Turkish, establishing the primacy of morphotactic rule compilation for handling unbounded affixation. For the Kipchak subgroup, Tyers and Washington~\citep{tyers2016finite} developed open-source finite-state transducers covering Kazakh, Tatar, and Kumyk using the HFST platform. Makhambetov et al.~\citep{makhambetov2013kazakh,makhambetov2015data} curated the first large-scale annotated Kazakh corpus and explored statistical data-driven morphological analyzers. More recently, Bekmanova et al.~\citep{bekmanova2021kazakh} formalized rule-based segmentation models combining stem dictionaries with semantic processing. However, prior finite-state systems focused almost exclusively on formal, canonical literary texts, failing to accommodate the widespread bilingual code-switching and foreign loanword affixes characteristic of modern digital communication.

\paragraph{Subword Tokenization Pathologies}
Modern deep learning architectures universally process raw text via statistical subword tokenization, notably Byte-Pair Encoding (BPE)~\citep{sennrich2016subword}, WordPiece~\citep{devlin2019bert}, and SentencePiece Unigram~\citep{kudo2018sentencepiece}. These algorithms iteratively merge frequent byte or character n-grams based on corpus frequency. While optimal for isolating or moderately inflected European languages, subword tokenizers introduce severe pathologies in agglutinative tongues. Because vocabulary frequency distributions in multilingual models (such as mBERT~\citep{devlin2019bert} and XLM-RoBERTa~\citep{conneau2020xlmr}) are dominated by high-resource languages, Kazakh affixation chains are fragmented into arbitrary subword tokens that rupture morphological boundaries. This subword token inflation obscures core semantic stems and disperses polarity markers across multiple tokens. While monolingual pretrained transformers (such as KazBERT~\citep{eraly2021kazbert} and KazRoBERTa~\citep{sagyndyk2025kazroberta}) mitigate out-of-vocabulary rates, they remain vulnerable to subword segmentation anomalies when encountering rare inflected forms or short informal expressions.

\subsection{Automated Fact Verification and Retrieval Pipelines}
Automated fact verification has advanced rapidly following the establishment of the Fact Extraction and VERification (FEVER) benchmark~\citep{thorne2018fever} and its structured successor FEVEROUS~\citep{aly2021feverous}. Canonical verification pipelines operate as a three-stage workflow: document retrieval, sentence selection, and natural language inference (NLI) stance classification predicting whether a claim is \textsc{Supported}, \textsc{Refuted}, or \textsc{Not Enough Info} (NEI).

\paragraph{Evidence Retrieval Paradigms}
Retrieval systems typically balance sparse term-frequency indexing and dense representation learning. BM25~\citep{robertson2009bm25} provides robust exact lexical matching but fails in morphologically rich languages when inflected query terms diverge from corpus surface forms. Conversely, dense retrievers such as Dense Passage Retrieval (DPR)~\citep{karpukhin2020dpr} and Contriever~\citep{izacard2021contriever} map claims and documents into continuous vector spaces via contrastive learning. However, dense encoders frequently struggle with exact numerical entities, dates, and morphologically nuanced negation in low-resource environments. Hybrid retrieval pipelines combining sparse and dense scoring via Reciprocal Rank Fusion~\citep{karpukhin2020dpr} have emerged as the state-of-the-art approach in open-domain question answering and knowledge-intensive NLP~\citep{lewis2020retrieval}.

\paragraph{Lexical Overlap Shortcuts and the NEI Trap}
A well-documented vulnerability in FEVER-style benchmarks is the presence of superficial dataset shortcuts~\citep{schuster2019fever}. Models frequently exploit spurious correlations: high unigram overlap between claim and evidence strongly correlates with \textsc{Supported}, whereas low unigram overlap indicates \textsc{NEI}. Schuster et al.~\citep{schuster2019fever} demonstrated that models trained on standard FEVER rely heavily on lexical matching heuristics, collapsing when evaluated on symmetric counterfactuals. In agglutinative languages, this vulnerability is catastrophic: modifying a single verbal negation allomorph (e.g., changing \textit{келді} ``arrived'' to \textit{келмеді} ``did not arrive'') completely reverses factual stance while retaining over 90\% of surface token overlap. The creation of robust benchmarks requiring morphological and evidential reasoning (such as SciFact~\citep{wadden2020fact}) has proven indispensable for evaluating true semantic comprehension.

\subsection{Multi-Dimensional Trust Modeling in Generative AI}
In contemporary information ecosystems, factual veracity and machine-generation likelihood represent orthogonal epistemic dimensions. A synthetic text generated by an LLM may be factually grounded and accurate (e.g., automated weather reports or verified translations), while authentic human prose may contain factual falsehoods, rumors, or unverified speculation. 

Traditional NLP systems isolate these tasks into independent silos: AI text detectors discard factual grounding, while fact-checkers operate under the assumption of human authorship. When foundation models hallucinate plausible citations, single-axis verification systems fail catastrophically. Recent investigations in uncertainty estimation and network calibration~\citep{guo2017calibration} emphasize the necessity of calibrated confidence scoring. Guo et al.~\citep{guo2026kazakh} established the conceptual foundation for multidimensional trust modeling in Central Asian languages. In this study, we build upon this foundation by formulating a continuous 2D Trust Matrix that maps factual veracity ($T_{\text{fact}}$) and generative authenticity ($T_{\text{gen}}$) into a joint Cartesian space, providing a unified mathematical framework for AI content auditing.
